\section{Diskussion}
\label{sec:Diskussion}

% Die relative Messabweichung wurde nach der Formel 
% \begin{equation}
%     \Delta x =   \frac{ \left| x_{exp} - x_{theo} \right|   }{x_{theo}}  
%     \label{eq:abweichung}
% \end{equation}
% bestimmt.


\subsection{Frequenzabhängigkeit der Leitungskonstanten}

Wie recht gut anhand der Plots zu sehen ist, sind die Leitungskonstanten $R,C,L$ größtenteils konstant für den Frequenzbereich $1-20 \mathrm{kHz}$. Bei $\qty{100}{\kilo\hertz}$ ist allerdings eine größere Abweichung zu erkennen, was zum Beispiel für den Widerstand $R$ durch den erst in diesem Bereich auftretenden Skin-Effekt erklärt werden kann. \\
